\begin{theorem}[Chebyshev’s sum inequality]
Let \(n\in\mathbb{N}\) and let  
\[
a_{1}\le a_{2}\le\cdots\le a_{n},\qquad 
b_{1}\le b_{2}\le\cdots\le b_{n}
\]
be real numbers. Then  
\begin{align}
n\sum_{i=1}^{n}a_{i}b_{i}\;\ge\;
\Bigl(\sum_{i=1}^{n}a_{i}\Bigr)\Bigl(\sum_{i=1}^{n}b_{i}\Bigr).
\end{align}
Equality holds \emph{iff} one of the following occurs
\begin{enumerate}
\item \(n=1\);
\item the sequence \((a_i)\) is constant, i.e. \(a_{1}=a_{2}= \dots =a_{n}\);
\item the sequence \((b_i)\) is constant, i.e. \(b_{1}=b_{2}= \dots =b_{n}\).
\end{enumerate}
\end{theorem}

\begin{proof}
\textbf{Trivial cases.}  
If \(n=1\) both sides equal \(a_{1}b_{1}\), so the inequality is an equality.  
If all \(a_i\) are equal (say \(a_i\equiv A\)), then
\[
n\sum_{i=1}^{n}a_i b_i = nA\sum_{i=1}^{n}b_i 
    = \Bigl(\sum_{i=1}^{n}a_i\Bigr)\Bigl(\sum_{i=1}^{n}b_i\Bigr),
\]
and similarly when all \(b_i\) are equal.  
Henceforth assume \(n\ge2\) and that neither \((a_i)\) nor \((b_i)\) is constant.

\medskip

\textbf{1. A non‑negative double sum.}  
Because the sequences are non‑decreasing, for any indices \(i<j\) we have
\[
a_i-a_j\le0,\qquad b_i-b_j\le0,
\]
and therefore
\begin{align}
\bigl(a_i-a_j\bigr)\bigl(b_i-b_j\bigr)\ge0
\qquad(1\le i<j\le n). \tag{1}
\end{align}
Summing \eqref{1} over all unordered pairs \((i,j)\) with \(i<j\) yields a non‑negative quantity
\begin{align}
S:=\sum_{1\le i<j\le n}(a_i-a_j)(b_i-b_j)\;\ge\;0 . \tag{2}
\end{align}

\medskip

\textbf{2. Expanding the double sum.}  
For each pair \((i,j)\),
\[
(a_i-a_j)(b_i-b_j)=a_i b_i + a_j b_j - a_i b_j - a_j b_i .
\]
Hence
\begin{align}
S
&=\sum_{i<j}\bigl(a_i b_i + a_j b_j - a_i b_j - a_j b_i\bigr) \\
&=\underbrace{\sum_{i<j}a_i b_i}_{A}
  \;+\;\underbrace{\sum_{i<j}a_j b_j}_{B}
  \;-\;\underbrace{\sum_{i<j}a_i b_j}_{C}
  \;-\;\underbrace{\sum_{i<j}a_j b_i}_{D}. \tag{3}
\end{align}

\textit{Counting the diagonal terms.}  
Each product \(a_k b_k\) appears once for every index that can be paired with a different index, i.e. \(n-1\) times. Consequently
\begin{align}
A+B=(n-1)\sum_{k=1}^{n} a_k b_k . \tag{4}
\end{align}

\textit{Counting the mixed terms.}  
Every mixed product \(a_i b_j\) with \(i\neq j\) occurs exactly once, either in \(C\) (when \(i<j\)) or in \(D\) (when \(j<i\)). Therefore
\begin{align}
C+D=\sum_{\substack{i,j=1\\ i\neq j}}^{n} a_i b_j
    =\Bigl(\sum_{i=1}^{n}a_i\Bigr)\Bigl(\sum_{j=1}^{n}b_j\Bigr)-\sum_{k=1}^{n} a_k b_k . \tag{5}
\end{align}

Substituting \eqref{4} and \eqref{5} into \eqref{3} gives
\begin{align}
S &= (n-1)\sum_{k=1}^{n} a_k b_k
     \;-\;\Bigl[\Bigl(\sum_{i=1}^{n}a_i\Bigr)\Bigl(\sum_{j=1}^{n}b_j\Bigr)
            -\sum_{k=1}^{n} a_k b_k\Bigr] \\
  &= n\sum_{k=1}^{n} a_k b_k
     \;-\;
     \Bigl(\sum_{i=1}^{n}a_i\Bigr)\Bigl(\sum_{j=1}^{n}b_j\Bigr). \tag{6}
\end{align}

\medskip

\textbf{3. Chebyshev’s inequality.}  
From \eqref{2} we have \(S\ge0\).  Using the identity \eqref{6},
\[
0\le S = n\sum_{i=1}^{n} a_i b_i 
          -\Bigl(\sum_{i=1}^{n}a_i\Bigr)\Bigl(\sum_{j=1}^{n}b_j\Bigr),
\]
which is exactly the claimed inequality.

\medskip

\textbf{4. Equality condition.}  
Equality \(S=0\) in \eqref{2} occurs precisely when each term \((a_i-a_j)(b_i-b_j)\) vanishes.  
For a fixed pair \(i<j\),
\[
(a_i-a_j)(b_i-b_j)=0 \iff a_i=a_j\ \text{or}\ b_i=b_j .
\]
If there exists a pair with \(a_i\neq a_j\), then we must have \(b_i=b_j\) for every such pair, forcing all \(b_k\) to be equal.  
Conversely, if some pair satisfies \(b_i\neq b_j\), then all \(a_k\) must be equal.  
Thus equality holds exactly in one of the mutually exclusive situations listed in the theorem statement:
\(n=1\), or the \(a\)-sequence is constant, or the \(b\)-sequence is constant.

\end{proof}